\chapter{Results/Evaluation}

\section{Software Testing}

\section{Interview-based Evaluation}

\subsection{Methodological Approach}
A qualitative interview-based evaluation was used, with semi-structured interviews.
This method was appropriate because the aim of the research was not only to gather feedback on the prototype, but also to investigate how participants made sense of new privacy controls, how they interpreted post-send revocation, and how they felt about the power balance between stakeholders. Weiss argues that qualitative interviewing is useful when the researcher needs more than brief or standardised answers, because it allows the interviewer to move beyond fixed questioning and obtain participants’ fuller accounts through further explanation and discussion (Weiss, 1994, p. 3). It was therefore well suited to this study, where participants needed to explain their own reasoning, expectations, and concerns in response to novel interface states. Weiss also notes that qualitative interviewing enables researchers to access participants’ “perceptions and reactions” that are only available through their own accounts (Weiss, 1994, p. 10).

Semi-structured interviews are well suited to this kind of exploratory HCI evaluation because they combine consistency across participants with enough flexibility to probe misunderstandings, emotional reactions, and unanticipated concerns.This is also consistent with Weiss’s view that, in qualitative interviewing, the interviewer should direct discussion toward the topics that matter to the study while also judging when a response requires elaboration (Weiss, 1994, pp. 8–9). In this study, the same scenario sequence was shown to all participants, while follow-up probes were used selectively depending on the participant's responses to explore their reasoning in more depth.

\subsection{Participants and Sampling}
The selection of participants were chosen using purposive sampling. This approach was appropriate because the study's goal was not to produce a statistically representative account of the general population. Instead, it seeks feedback from people whose daily experiences qualified them to evaluate a prototype for messaging-based photo sharing. Weiss1994 argues that, when a study is concerned with the experiences of people who share a common characteristic, the researcher needs “a sample of people who together can represent the population of concern,” and that with small qualitative samples it is often preferable to select respondents purposively in order to capture important variation rather than rely on probability sampling \cite{weiss_1994}.

The study specifically targeted adults aged 18-24, as this group is especially relevant to the context of the prototype and to contemporary messaging-based photo-sharing practices. 
In the UK, 18-24-year-olds remained the most active online in 2025, while 16-24-year-olds used the widest range of online communication platforms of any age group, indicating particularly high engagement with digital communication services \cite{OfcomOnlineNation2025,OfcomAdultsMediaUse2025}. This is directly relevant to the current study because online communication platforms include services used for messaging and calling, and WhatsApp has been identified as both the most popular app for those activities and the highest-reaching messaging and calls service across age groups \cite{OfcomOnlineNation2024,OfcomAdultsMediaUse2025}. Although some usage data is reported for the broader 16-24 category, the present study focuses on the adult subset of that group for practical and ethical reasons. Taken together, these findings support the use of an 18-24 sample for evaluating a prototype concerned with image-sharing permissions and consent revocation in messaging environments.

A target of 8-10 participants was set for this study. This number was considered appropriate for a focused qualitative evaluation with a relatively specific participant group and a narrow research aim. Rather than seeking statistical outcomes, the study aimed to identify the recurring patterns in how participants understood the prototype. This is supported by methodological literature suggesting that smaller interview samples can be sufficient in focused studies. Hennink, Kaiser and Marconi found that code saturation was reached at nine interviews, while Guest, Bunce and Johnson found that thematic saturation occured within the first twelwe interviws, with basic elements of higher-level themes were already present by six interviews \cite{Hennink_2017,Guest_Bunce_Johnson_2006}. More broadly, Malterud, Siersma and Guassora argue that qualitative sample size should be guided by information power, such that studies with a narrow aim and specific sample may require fewer participants \cite{malterud_2016}. In this case, 8--10 participants was judged sufficient to generate meaningful evaluative insight while remaining feasible within the scope of the dissertation. 



